\chapter{堆積学}
\section{背景：失われた流れを岩石から推定する}
\label{theory:sediment:background}
地層には、形成時の水流そのものは残らない。残るのは粒径、形、鉱物組成、層の幾何、生物の痕跡である。堆積学は、粒子の運搬と沈積という順問題を使い、現在観察できる地層から過去の環境を推定する逆問題である。同じ砂岩でも河川、砂丘、浅海など複数の候補がある。まず観測した特徴を書き、次に成因の仮説を立てる。観測と解釈を一語に押し込めない。\cite{sed_mit_environment}

\section{堆積相、組合せ、累重}
\label{theory:sediment:facies}
堆積相とは、岩相、粒径、堆積構造、化石などの組合せによって周囲と区別する地層の部分である。「細粒砂岩、斜交葉理、植物片」という記載と「河川堆積物」という解釈を分ける。相の組合せは堆積相アソシエーション、上下の配列は堆積相サクセションである。堆積シーケンスは単に相を集めたものではなく、不整合又はそれに対比される整合面を境界とする成因的に関連した地層群をいう。尺度と境界条件が異なる。\cite{sed_mit_environment,sed_walther_oregon}

\section{ワルターの法則：空間的隣接が時間的累重になる}
\label{theory:sediment:walther}
浅い所の砂と深い所の泥が側方に隣接して堆積している状況を考える。環境帯が移動すると、一地点では泥の上へ砂が重なる。連続した整合的・漸移的な累重において、上下の相はもともと側方で隣接する環境に対応するというのがワルターの法則の骨子である。これは全ての上下接触へ使える規則ではない。不整合や大きな侵食面を挟めば、長い時間後の別の環境が重なるため、上下関係だけから同時の側方関係を復元できない。\cite{sed_walther_oregon}

\section{粒子を動かす力と沈降}
\label{theory:sediment:transport}
水流は底面せん断によって粒子を動かし、乱流は沈降する粒子を再び上へ運ぶ。粒子は底面に沿う掃流と、水中の浮遊に分けて考えられる。球形の孤立粒子が低レイノルズ数で沈む理想条件では、重力と浮力の差を粘性抵抗と釣り合わせると
\[
 w_s=\frac{(\rho_s-\rho_f)gd^2}{18\eta}
\]
となる。\(w_s\) は終端沈降速度、\(d\) は直径、\(\rho_s,\rho_f\) は粒子と流体の密度、\(\eta\) は粘性係数である。粒径が大きいほど速く沈むという見通しを与えるが、粗い砂礫、非球形、高濃度の群れへこの式を直接代入しない。堆積構造は沈降だけでなく、底面の侵食と粒子供給の時間変化にも依存する。

\section{河川の二つの端成分と洪水}
\label{theory:sediment:river}
蛇行河川では湾曲した流路が側方へ移動し、内岸のポイントバーに砂がたまり、流路外の氾濫原には泥がたまりやすい。網状河川では複数の流路が砂礫州を挟み、粗粒負荷や流量変動の大きい条件が典型となる。しかし二つは排他的な分類ではなく、中間形態もある。勾配だけで流路を決めず、河岸の強さ、供給量、粒度を合わせて見る。\cite{sed_mit_environment}

\section{正級化と逆級化を区別する}
\label{theory:sediment:grading}
一枚の層の中で上へ細粒化するのが正級化、上へ粗粒化するのが逆級化である。流れの減衰と粗粒子の先行沈降は正級化の典型機構である。洪水の立上がりでは運搬可能粒径が時間とともに増し、氾濫原へ届く砂が上方粗粒化することがある。だが逆級化は土石流などにも見られ、河川だけの指標ではない。一枚のイベント層内の級化と、複数の相が重なる上方粗粒化の累重を区別する。

\section{デルタ、海進、海退と河口内湾}
\label{theory:sediment:delta}
デルタは河川が海や湖へ土砂を供給してできる河口の堆積体である。前進すると海側の細粒層の上へ浅い粗粒層が重なる。海退は海岸線が海側へ移ることであり、海水準の低下だけを意味しない。相対海水準が上がっても、土砂供給が新しく生じる堆積空間を埋める速さを上回れば前進する通常海退となる。相対海水準低下による強制海退と分ける。海進で河口や谷が沈水して形成されるエスチュアリーでは、河川と潮汐の双方が働く。\cite{sed_walther_oregon}

\section{移動するリップルと斜交層理}
\label{theory:sediment:cross}
流れが砂の床形を移動させると、粒子は緩い上流側を運ばれ、下流側の斜面へたまる。その斜面が保存されると斜交葉理又は斜交層理となる。小さな床形のリップルも大きなデューンも同じ幾何の考え方を使う。平板型とトラフ型の違いは床形の三次元性や切込みの形に関わり、別々の流体力学法則が必要という意味ではない。古流向は層の傾斜を補正し、複数断面から求める。\cite{sed_houston_structures}

\section{ヘリンボーン：往復流と断面の錯覚}
\label{theory:sediment:herringbone}
逆向きに傾く斜交層の組は、潮流のように向きが反転する流れを考える手掛かりとなる。ヘリンボーンは蛇行河川の一方向流の代表構造ではない。ただし一断面で逆向きに見えることだけでは、真の180度反転を保証しない。三次元の床形を斜めに切った断面でも似た形が見える。再活動面や泥の薄い覆い、古流向の分布を併用する。\cite{sed_mit_environment}

\section{ハンモック状斜交層理と暴浪}
\label{theory:sediment:hcs}
ゆるい凸部と凹部へ沿う低角度の層理がハンモック状斜交層理の特徴である。典型的にはストーム時の波浪による振動流、又は一方向成分を含む複合流を考える。堆積構造の形の名前と成因の名前を同一視しない。MITの講義では広い意味の同構造が多成因であることと、反砂堆に関わる例も扱う。したがって試験で「一方向流の代表として最も不適切」と選ぶことは、「一方向成分のある全ての条件で絶対に形成されない」という断定ではない。\cite{sed_mit_cross}

\section{混濁流とタービダイト}
\label{theory:sediment:turbidite}
混濁流は浮遊粒子によって周囲の水より密度が高くなった流れが、重力で斜面に沿って進む現象である。乱流が粒子の沈降に対抗する。沈積した層をタービダイトと呼び、流れが弱まると正級化が生じやすい。バウマシーケンスは級化砂、平行葉理、リップル葉理、細粒部などの一イベント内の典型配列である。全区分が毎回そろうわけではなく、侵食や堆積しない区分の省略がある。海底地すべりは発生源の一つであって、全ての混濁流の必要条件ではない。\cite{sed_mit_environment}

\section{未固結変形、荷重痕、火炎、脱水}
\label{theory:sediment:soft}
まだ固まらない泥の上へ砂が急速にたまると、重い砂が下へ突出する荷重痕と、泥が上へ伸びる火炎構造ができることがある。堆積後の未固結段階の変形であるため、粒子をたまらせる過程の葉理・級化・リップル・流痕と区別する。脱水による皿状構造も、層の記載だけでなく間隙水の逃げ方を考える対象である。液状化を伴う場合、有効応力 \(\sigma'=\sigma-u\) の低下が変形しやすさを説明する。\(u\) は間隙水圧である。火炎があるだけで地震起源を断定せず、過荷重などの候補と比較する。\cite{sed_soft_deformation,sed_houston_structures}

\section{地層の上下判定}
\label{theory:sediment:wayup}
現在の鉛直方向と形成時の上は、褶曲や転倒で一致しない。荷重痕は砂が下の泥へ突き出る向き、皿状構造は一般的な上に凹の形、流痕は侵食された面とその上の層が埋めた凸型の底面を手掛かりとする。生痕も巣穴の開口や分岐・埋積の幾何が分かる場合に上下面を示すが、名称だけでは全ての生痕に向きを付けられない。節理面の羽毛状構造は割れ目が広がった方向を示し、堆積時の上下を示すものではない。\cite{sed_houston_structures,sed_plumose_michigan}

\section{環境指標と分析装置の観測量}
\label{theory:sediment:proxy}
黄鉄鉱 \(\mathrm{FeS_2}\) の形成には鉄・硫黄の供給と還元状態が関わる。海成粘土を酸化・乾燥して懸濁すると硫酸塩の生成などが電気伝導度へ影響するため、古環境推定に利用される。だが作成手順、濃度、温度、後の溶脱も値を変えるため、海成と淡水成の万能な閾値を設けない。\cite{sed_gsj_ec2011}
蛍光X線分析は特性X線のエネルギーと強度から元素を調べる方法であり、内部の堆積構造を直接撮影する方法ではない。SEMは粒子表面の微細形状を観察する。運搬で生じた摩耗や衝突の痕跡は風成作用を検討する材料となるが、溶食や再運搬でも変わる。装置の名前、観測量、環境解釈という三段を区別する。\cite{sed_xrf_jaea,sed_sem_shimane}

\section{津波の地質記録：形成と保存は別である}
\label{theory:sediment:tsunami}
津波は海底や海岸の粒子を陸側へ運び、戻り流れは再び海側へ運ぶ。複数波、再侵食、滞留時間の差によって、塊状、葉理、級化、複数ユニットなどが生じる。一方、近年のUSGSデータ集では構造が稀な例も多い。したがって「単純が常に多い」「複雑が常に多い」のどちらにも一般化しない。\cite{sed_tsunami_fujino,sed_tsunami_usgs}
保存は形成後の侵食と再攪乱に依存する。静かな湿地や内湾の細粒背景層に異質な砂層が挟まれば調べやすいが、砂層だけで津波と認定しない。微化石、供給域、広がり、複数の構造、年代を合わせる。生物遺骸は供給源の水深を検討する手掛かりであり、混合・再堆積の影響を受ける。岩礁や礫の多い場所も記録を持ち得るが、識別と継続保存が難しい。\cite{sed_tsunami_2002,sed_tsunami_gsj_survey}

\section{典型的な問われ方}
過去問では、構造と流れの対応、初生と二次の区別、上下判定、分析手法の観測量、相・累重・シーケンスの名称の入替えが問われる。正答番号は同じテーマでも年度で変わる。全肢を条件付きで読むことに加え、「最も不適切」と科学的な全称命題を分ける。頻度は全問の分類確認後に集計する。

\section{過去問と個別の解法}
\section{流れ・波浪と斜交層理}

\subsection{平成30年度 問25}
\label{q:h30-reg-applied-25}

\noindent\textbf{原典：} sources/H30/17\_applied\_science.pdf、PDF物理8頁。

\textbf{公式正答：} 2（single）。

\par\textbf{検討未確定：} 海底地すべりを全タービダイトの起源、全Bouma区分を必須とする読みは過度である。HCSにも多成因がある。公式/独立2を保持して典型対応の範囲を部分校正に残す。

\paragraph{問題の要旨} 流れ・波浪と堆積構造を対応させ、蛇行河川と往復する潮流の構造を区別する。

\Needspace{4\baselineskip}\paragraph{確認した条件・記述（独自要約）}\mbox{}\par

\noindent\textbf{判断方向：} 最も不適切・誤りを選ぶ\par

\noindent\textbf{原典照合：} 15原典画像で18問の問題幹、全5肢、否定、専門語、化学式、候補順を確認。\par

\noindent\textbf{適用条件：} 流れと構造の典型対応は一意成因を保証しない。用語、時間間隙、保存、試料処理の条件を個別に示す。\par

\paragraph{問われる核心} 流れ・波浪と堆積構造を対応させ、蛇行河川と往復する潮流の構造を区別する。

\paragraph{必要知識} 床形の移動と斜交、往復流、振動・複合流、密度流、流痕、典型成因の適用条件。

\paragraph{考え方・正答を選ぶ理由（査読待ち初稿）} 床形の前進は斜交層理を保存し、流れの侵食は流痕を作る。混濁流の減衰は一般に級化したタービダイトを残す。ハンモック状斜交層理では暴浪時の振動流又は複合流が代表となる。ヘリンボーンを蛇行河川の主な構造に割り当てる対応が最も不適切であり、本問の番号は2である。潮流の反転を典型とするが、一断面の逆傾斜だけで真の180度反転を断定しない。原典のタービダイト説明も全例の起源を限定せず読む。

\Needspace{4\baselineskip}\paragraph{選択肢の検討（査読待ち初稿）}\mbox{}\par

\noindent\textbf{1：適切（条件を含めて検討）} 床形の移動で前面へ層が付き斜交層理となる。移動する床形の下流側斜面への堆積が内部の斜交を保存するという説明に対応する。\par

\noindent\textbf{2：誤り・最も不適切} 蛇行河川の代表としてヘリンボーンを挙げる。相反する流向で傾斜の異なる層が重なる構造は潮流を典型とし、蛇行の一方向流の主な構造という対応が違う。\par

\noindent\textbf{3：適切（条件を含めて検討）} 暴浪の波浪でハンモック状斜交層理が形成する。ストームの振動流・複合流が代表的な形成条件である。形の名称から全例の単一成因を断定しない。\par

\noindent\textbf{4：適切（条件を含めて検討）} 陸棚斜面の海底地すべりに起源を持つ混濁流の堆積物がタービダイトで級化・Boumaが特徴。海底地すべりは混濁流の発生源の一つである。級化とBoumaは典型だが全ての混濁流の起源や全区分の存在を限定する定義ではない。\par

\noindent\textbf{5：適切（条件を含めて検討）} 流れの小さな削込みを流痕とする。流れが底面を侵食した痕跡であり、充填された型が上位層底面に保存される場合もある。\par

\paragraph{再利用する考え方} 観察できる構造と成因・環境の解釈を区別する。代表モデルに適用条件を添え、単一の形や測定値から唯一の環境を断定しない。

\paragraph{共通理論} \hyperref[theory:sediment:cross]{第\ref*{theory:sediment:cross}節}, \hyperref[theory:sediment:herringbone]{第\ref*{theory:sediment:herringbone}節}, \hyperref[theory:sediment:hcs]{第\ref*{theory:sediment:hcs}節}, \hyperref[theory:sediment:turbidite]{第\ref*{theory:sediment:turbidite}節}

\par 参考：\cite{sed_mit_environment,sed_mit_cross,sed_houston_structures}。

\paragraph{関連問題} \hyperref[q:r04-reg-applied-28]{R04-REG-APPLIED-28}

\par 難易度 3（校準初稿）、優先度未評価。

\subsection{令和4年度 問28}
\label{q:r04-reg-applied-28}

\noindent\textbf{原典：} sources/R04/17\_applied\_science.pdf、PDF物理10頁。

\textbf{公式正答：} 1（single）。

\paragraph{問題の要旨} 流れ・波浪と堆積構造を対応させ、蛇行河川と往復する潮流の構造を区別する。

\Needspace{4\baselineskip}\paragraph{確認した条件・記述（独自要約）}\mbox{}\par

\noindent\textbf{判断方向：} 最も不適切・誤りを選ぶ\par

\noindent\textbf{原典照合：} 15原典画像で18問の問題幹、全5肢、否定、専門語、化学式、候補順を確認。\par

\noindent\textbf{適用条件：} 流れと構造の典型対応は一意成因を保証しない。用語、時間間隙、保存、試料処理の条件を個別に示す。\par

\paragraph{問われる核心} 流れ・波浪と堆積構造を対応させ、蛇行河川と往復する潮流の構造を区別する。

\paragraph{必要知識} 床形の移動と斜交、往復流、振動・複合流、密度流、流痕、典型成因の適用条件。

\paragraph{考え方・正答を選ぶ理由（査読待ち初稿）} 床形の前進は斜交層理を保存し、流れの侵食は流痕を作る。混濁流の減衰は一般に級化したタービダイトを残す。ハンモック状斜交層理では暴浪時の振動流又は複合流が代表となる。ヘリンボーンを蛇行河川の主な構造に割り当てる対応が最も不適切であり、本問の番号は1である。潮流の反転を典型とするが、一断面の逆傾斜だけで真の180度反転を断定しない。原典のタービダイト説明も全例の起源を限定せず読む。

\Needspace{4\baselineskip}\paragraph{選択肢の検討（査読待ち初稿）}\mbox{}\par

\noindent\textbf{1：誤り・最も不適切} 蛇行河川の代表としてヘリンボーンを挙げる。相反する流向で傾斜の異なる層が重なる構造は潮流を典型とし、蛇行の一方向流の主な構造という対応が違う。\par

\noindent\textbf{2：適切（条件を含めて検討）} 暴浪の波浪でハンモック状斜交層理が形成する。ストームの振動流・複合流が代表的な形成条件である。形の名称から全例の単一成因を断定しない。\par

\noindent\textbf{3：適切（条件を含めて検討）} 床形の移動で前面へ層が付き斜交層理となる。移動する床形の下流側斜面への堆積が内部の斜交を保存するという説明に対応する。\par

\noindent\textbf{4：適切（条件を含めて検討）} 流れの小さな削込みを流痕とする。流れが底面を侵食した痕跡であり、充填された型が上位層底面に保存される場合もある。\par

\noindent\textbf{5：適切（条件を含めて検討）} 混濁流がタービダイトを作り一般に級化する。密度差で流れ、乱流が粒子を支持し、流れの減衰と沈降で級化しやすい。全例で全Bouma区分がそろうとは限らない。\par

\paragraph{再利用する考え方} 観察できる構造と成因・環境の解釈を区別する。代表モデルに適用条件を添え、単一の形や測定値から唯一の環境を断定しない。

\paragraph{共通理論} \hyperref[theory:sediment:cross]{第\ref*{theory:sediment:cross}節}, \hyperref[theory:sediment:herringbone]{第\ref*{theory:sediment:herringbone}節}, \hyperref[theory:sediment:hcs]{第\ref*{theory:sediment:hcs}節}, \hyperref[theory:sediment:turbidite]{第\ref*{theory:sediment:turbidite}節}

\par 参考：\cite{sed_mit_environment,sed_mit_cross,sed_houston_structures}。

\paragraph{関連問題} \hyperref[q:h30-reg-applied-25]{H30-REG-APPLIED-25}, \hyperref[q:r07-reg-applied-27]{R07-REG-APPLIED-27}

\par 難易度 3（校準初稿）、優先度未評価。

\subsection{令和5年度 問31}
\label{q:r05-reg-applied-31}

\noindent\textbf{原典：} sources/R05/17\_applied\_science.pdf、PDF物理11頁。

\textbf{公式正答：} 2（single）。

\par\textbf{検討未確定：} 広義のHCSには多成因があり、一方向流で絶対に不可能という全称は未支持である。初回/公式2は典型モデルの比較として保持し、用語範囲をAstraへ残す。

\paragraph{問題の要旨} 一方向流で形成し得る層理と、波浪の振動・複合流が典型となる構造を比較する。

\Needspace{4\baselineskip}\paragraph{確認した条件・記述（独自要約）}\mbox{}\par

\noindent\textbf{判断方向：} 最も不適切・誤りを選ぶ\par

\noindent\textbf{原典照合：} 15原典画像で18問の問題幹、全5肢、否定、専門語、化学式、候補順を確認。\par

\noindent\textbf{適用条件：} 流れと構造の典型対応は一意成因を保証しない。用語、時間間隙、保存、試料処理の条件を個別に示す。\par

\paragraph{問われる核心} 一方向流で形成し得る層理と、波浪の振動・複合流が典型となる構造を比較する。

\paragraph{必要知識} 一方向の床形と層理、振動流と複合流、形の分類と成因の非一意性。

\paragraph{考え方・正答を選ぶ理由（査読待ち初稿）} 四つは一方向の掃流や床形の移動で形成し得るのに対し、ハンモック状斜交層理は波浪の振動流・複合流の代表なので2を選ぶ。ただしMITは広義のHCSに反砂堆を伴う例を含む多成因を挙げる。「最も不適切」という試験の代表比較を、全ての一方向流条件で絶対に形成されないという証明へ置き換えない。

\Needspace{4\baselineskip}\paragraph{選択肢の検討（査読待ち初稿）}\mbox{}\par

\noindent\textbf{1：適切（条件を含めて検討）} トラフ型斜交層理。三次元の床形の移動・侵食と堆積により一方向流で形成し得る。\par

\noindent\textbf{2：誤り・最も不適切} ハンモック状斜交層理。典型は暴浪の振動・複合流であり、本問で一方向流の代表として最も不適切である。ただし広義のHCSには複数成因が報告される。\par

\noindent\textbf{3：適切（条件を含めて検討）} フォーセット層理。床形の下流側斜面への堆積などにより一方向流で形成し得る。\par

\noindent\textbf{4：適切（条件を含めて検討）} 平行層理。床面に沿うシート状の運搬・堆積などで形成し得る。一方向流だけの固有指標という意味ではない。\par

\noindent\textbf{5：適切（条件を含めて検討）} 平板型斜交層理。二次元に近い床形の移動により一方向流で形成し得る。\par

\paragraph{再利用する考え方} 観察できる構造と成因・環境の解釈を区別する。代表モデルに適用条件を添え、単一の形や測定値から唯一の環境を断定しない。

\paragraph{共通理論} \hyperref[theory:sediment:cross]{第\ref*{theory:sediment:cross}節}, \hyperref[theory:sediment:hcs]{第\ref*{theory:sediment:hcs}節}

\par 参考：\cite{sed_mit_cross,sed_houston_structures}。

\par 難易度 3（校準初稿）、優先度未評価。

\subsection{令和7年度 問27}
\label{q:r07-reg-applied-27}

\noindent\textbf{原典：} sources/R07/17\_applied\_science.pdf、PDF物理10頁。

\textbf{公式正答：} 3（single）。

\paragraph{問題の要旨} 流れ・波浪と堆積構造を対応させ、蛇行河川と往復する潮流の構造を区別する。

\Needspace{4\baselineskip}\paragraph{確認した条件・記述（独自要約）}\mbox{}\par

\noindent\textbf{判断方向：} 最も不適切・誤りを選ぶ\par

\noindent\textbf{原典照合：} 15原典画像で18問の問題幹、全5肢、否定、専門語、化学式、候補順を確認。\par

\noindent\textbf{適用条件：} 流れと構造の典型対応は一意成因を保証しない。用語、時間間隙、保存、試料処理の条件を個別に示す。\par

\paragraph{問われる核心} 流れ・波浪と堆積構造を対応させ、蛇行河川と往復する潮流の構造を区別する。

\paragraph{必要知識} 床形の移動と斜交、往復流、振動・複合流、密度流、流痕、典型成因の適用条件。

\paragraph{考え方・正答を選ぶ理由（査読待ち初稿）} 床形の前進は斜交層理を保存し、流れの侵食は流痕を作る。混濁流の減衰は一般に級化したタービダイトを残す。ハンモック状斜交層理では暴浪時の振動流又は複合流が代表となる。ヘリンボーンを蛇行河川の主な構造に割り当てる対応が最も不適切であり、本問の番号は3である。潮流の反転を典型とするが、一断面の逆傾斜だけで真の180度反転を断定しない。原典のタービダイト説明も全例の起源を限定せず読む。

\Needspace{4\baselineskip}\paragraph{選択肢の検討（査読待ち初稿）}\mbox{}\par

\noindent\textbf{1：適切（条件を含めて検討）} 流れの小さな削込みを流痕とする。流れが底面を侵食した痕跡であり、充填された型が上位層底面に保存される場合もある。\par

\noindent\textbf{2：適切（条件を含めて検討）} 床形の移動で前面へ層が付き斜交層理となる。移動する床形の下流側斜面への堆積が内部の斜交を保存するという説明に対応する。\par

\noindent\textbf{3：誤り・最も不適切} 蛇行河川の代表としてヘリンボーンを挙げる。相反する流向で傾斜の異なる層が重なる構造は潮流を典型とし、蛇行の一方向流の主な構造という対応が違う。\par

\noindent\textbf{4：適切（条件を含めて検討）} 混濁流がタービダイトを作り一般に級化する。密度差で流れ、乱流が粒子を支持し、流れの減衰と沈降で級化しやすい。全例で全Bouma区分がそろうとは限らない。\par

\noindent\textbf{5：適切（条件を含めて検討）} 暴浪の波浪でハンモック状斜交層理が形成する。ストームの振動流・複合流が代表的な形成条件である。形の名称から全例の単一成因を断定しない。\par

\paragraph{再利用する考え方} 観察できる構造と成因・環境の解釈を区別する。代表モデルに適用条件を添え、単一の形や測定値から唯一の環境を断定しない。

\paragraph{共通理論} \hyperref[theory:sediment:cross]{第\ref*{theory:sediment:cross}節}, \hyperref[theory:sediment:herringbone]{第\ref*{theory:sediment:herringbone}節}, \hyperref[theory:sediment:hcs]{第\ref*{theory:sediment:hcs}節}, \hyperref[theory:sediment:turbidite]{第\ref*{theory:sediment:turbidite}節}

\par 参考：\cite{sed_mit_environment,sed_mit_cross,sed_houston_structures}。

\paragraph{関連問題} \hyperref[q:r04-reg-applied-28]{R04-REG-APPLIED-28}

\par 難易度 3（校準初稿）、優先度未評価。

\section{堆積相と河川・デルタ・河口システム}

\subsection{平成23年度 問25}
\label{q:h23-reg-applied-25}

\noindent\textbf{原典：} sources/H23/17\_applied\_science.pdf、PDF物理8頁。

\textbf{公式正答：} 5（single）。

\par\textbf{検討未確定：} デルタを海退期と呼ぶ範囲、氾濫原砂の逆級化の代表性は要限定である。逆級化は他成因にもあり河川に一意ではない。全河川・全デルタの排他的な診断とはしない。独立初回と公式の蛇行/網状混同肢を保持する。

\paragraph{問題の要旨} 河川と河口の堆積システムを比較し、蛇行と網状の典型条件を区別する。

\Needspace{4\baselineskip}\paragraph{確認した条件・記述（独自要約）}\mbox{}\par

\noindent\textbf{判断方向：} 最も不適切・誤りを選ぶ\par

\noindent\textbf{原典照合：} 15原典画像で18問の問題幹、全5肢、否定、専門語、化学式、候補順を確認。\par

\noindent\textbf{適用条件：} 流れと構造の典型対応は一意成因を保証しない。用語、時間間隙、保存、試料処理の条件を個別に示す。\par

\paragraph{問われる核心} 河川と河口の堆積システムを比較し、蛇行と網状の典型条件を区別する。

\paragraph{必要知識} 河川端成分、氾濫原とポイントバー、洪水時の級化、デルタ前進、相対海水準と供給、河口沈水。

\paragraph{考え方・正答を選ぶ理由（査読待ち初稿）} 五つは河川の端成分と河口堆積システムの対応を問う。蛇行の典型を、急勾配・大きな流量変動・粗粒子の大きな負荷で説明する肢を除く。この組合せは網状の典型に対応する。ただし実際の流路型には河岸の強度や供給も関わる。デルタは前進する堆積体で、海退と相対海水準低下は同義ではない。逆級化も洪水立上がりで形成し得るが、非一意の指標である。原典の典型モデルと、全称命題としての限界を分けて解く。本問の該当番号は5である。

\Needspace{4\baselineskip}\paragraph{選択肢の検討（査読待ち初稿）}\mbox{}\par

\noindent\textbf{1：適切（条件を含めて検討）} 河川を蛇行・網状の二端成分で考える。端成分モデルとして妥当である。中間形態や両特徴を持つ河川を除外する絶対二分ではない。\par

\noindent\textbf{2：適切（条件を含めて検討）} 海湖側へ突出するデルタを海退期のシステムとする。デルタの前進が海退を示す場合の説明である。海退を海水準低下と同義にせず、デルタの存在だけで全時期を決めない。\par

\noindent\textbf{3：適切（条件を含めて検討）} 海進で河口が沈水する内湾をエスチュアリーとする。沈水谷を河川と潮汐が埋める河口内湾というモデルに対応する。海進する全河口の唯一の形ではない。\par

\noindent\textbf{4：適切（条件を含めて検討）} 氾濫原の砂岩に見られる逆級化を河川の特徴とする。洪水の増水過程では砂が上方粗粒化する場合がある。ただし逆級化は河川に固有ではなく、全氾濫原砂がそうなるという意味ではない。\par

\noindent\textbf{5：誤り・最も不適切} 蛇行河川を流量変動・勾配・粗粒負荷の大きい河川とする。これは網状河川の典型条件との混同である。蛇行は流路の曲率と側方移動、ポイントバーと泥質氾濫原で捉える。\par

\paragraph{再利用する考え方} 観察できる構造と成因・環境の解釈を区別する。代表モデルに適用条件を添え、単一の形や測定値から唯一の環境を断定しない。

\paragraph{共通理論} \hyperref[theory:sediment:river]{第\ref*{theory:sediment:river}節}, \hyperref[theory:sediment:grading]{第\ref*{theory:sediment:grading}節}, \hyperref[theory:sediment:delta]{第\ref*{theory:sediment:delta}節}

\par 参考：\cite{sed_mit_environment,sed_walther_oregon}。

\paragraph{関連問題} \hyperref[q:h29-reg-applied-25]{H29-REG-APPLIED-25}, \hyperref[q:r06-reg-applied-31]{R06-REG-APPLIED-31}

\par 難易度 3（校準初稿）、優先度未評価。

\subsection{平成29年度 問25}
\label{q:h29-reg-applied-25}

\noindent\textbf{原典：} sources/H29/17\_applied\_science.pdf、PDF物理8頁。

\textbf{公式正答：} 5（single）。

\par\textbf{検討未確定：} デルタを海退期と呼ぶ範囲、氾濫原砂の逆級化の代表性は要限定である。逆級化は他成因にもあり河川に一意ではない。全河川・全デルタの排他的な診断とはしない。独立初回と公式の蛇行/網状混同肢を保持する。

\paragraph{問題の要旨} 河川と河口の堆積システムを比較し、蛇行と網状の典型条件を区別する。

\Needspace{4\baselineskip}\paragraph{確認した条件・記述（独自要約）}\mbox{}\par

\noindent\textbf{判断方向：} 最も不適切・誤りを選ぶ\par

\noindent\textbf{原典照合：} 15原典画像で18問の問題幹、全5肢、否定、専門語、化学式、候補順を確認。\par

\noindent\textbf{適用条件：} 流れと構造の典型対応は一意成因を保証しない。用語、時間間隙、保存、試料処理の条件を個別に示す。\par

\paragraph{問われる核心} 河川と河口の堆積システムを比較し、蛇行と網状の典型条件を区別する。

\paragraph{必要知識} 河川端成分、氾濫原とポイントバー、洪水時の級化、デルタ前進、相対海水準と供給、河口沈水。

\paragraph{考え方・正答を選ぶ理由（査読待ち初稿）} 五つは河川の端成分と河口堆積システムの対応を問う。蛇行の典型を、急勾配・大きな流量変動・粗粒子の大きな負荷で説明する肢を除く。この組合せは網状の典型に対応する。ただし実際の流路型には河岸の強度や供給も関わる。デルタは前進する堆積体で、海退と相対海水準低下は同義ではない。逆級化も洪水立上がりで形成し得るが、非一意の指標である。原典の典型モデルと、全称命題としての限界を分けて解く。本問の該当番号は5である。

\Needspace{4\baselineskip}\paragraph{選択肢の検討（査読待ち初稿）}\mbox{}\par

\noindent\textbf{1：適切（条件を含めて検討）} 河川を蛇行・網状の二端成分で考える。端成分モデルとして妥当である。中間形態や両特徴を持つ河川を除外する絶対二分ではない。\par

\noindent\textbf{2：適切（条件を含めて検討）} 海湖側へ突出するデルタを海退期のシステムとする。デルタの前進が海退を示す場合の説明である。海退を海水準低下と同義にせず、デルタの存在だけで全時期を決めない。\par

\noindent\textbf{3：適切（条件を含めて検討）} 海進で河口が沈水する内湾をエスチュアリーとする。沈水谷を河川と潮汐が埋める河口内湾というモデルに対応する。海進する全河口の唯一の形ではない。\par

\noindent\textbf{4：適切（条件を含めて検討）} 氾濫原の砂岩に見られる逆級化を河川の特徴とする。洪水の増水過程では砂が上方粗粒化する場合がある。ただし逆級化は河川に固有ではなく、全氾濫原砂がそうなるという意味ではない。\par

\noindent\textbf{5：誤り・最も不適切} 蛇行河川を流量変動・勾配・粗粒負荷の大きい河川とする。これは網状河川の典型条件との混同である。蛇行は流路の曲率と側方移動、ポイントバーと泥質氾濫原で捉える。\par

\paragraph{再利用する考え方} 観察できる構造と成因・環境の解釈を区別する。代表モデルに適用条件を添え、単一の形や測定値から唯一の環境を断定しない。

\paragraph{共通理論} \hyperref[theory:sediment:river]{第\ref*{theory:sediment:river}節}, \hyperref[theory:sediment:grading]{第\ref*{theory:sediment:grading}節}, \hyperref[theory:sediment:delta]{第\ref*{theory:sediment:delta}節}

\par 参考：\cite{sed_mit_environment,sed_walther_oregon}。

\paragraph{関連問題} \hyperref[q:h23-reg-applied-25]{H23-REG-APPLIED-25}, \hyperref[q:r05-reg-applied-30]{R05-REG-APPLIED-30}

\par 難易度 3（校準初稿）、優先度未評価。

\subsection{令和5年度 問30}
\label{q:r05-reg-applied-30}

\noindent\textbf{原典：} sources/R05/17\_applied\_science.pdf、PDF物理11頁。

\textbf{公式正答：} 3（single）。

\par\textbf{検討未確定：} デルタを海退期と呼ぶ範囲、氾濫原砂の逆級化の代表性は要限定である。逆級化は他成因にもあり河川に一意ではない。全河川・全デルタの排他的な診断とはしない。独立初回と公式の蛇行/網状混同肢を保持する。

\paragraph{問題の要旨} 河川と河口の堆積システムを比較し、蛇行と網状の典型条件を区別する。

\Needspace{4\baselineskip}\paragraph{確認した条件・記述（独自要約）}\mbox{}\par

\noindent\textbf{判断方向：} 最も不適切・誤りを選ぶ\par

\noindent\textbf{原典照合：} 15原典画像で18問の問題幹、全5肢、否定、専門語、化学式、候補順を確認。\par

\noindent\textbf{適用条件：} 流れと構造の典型対応は一意成因を保証しない。用語、時間間隙、保存、試料処理の条件を個別に示す。\par

\paragraph{問われる核心} 河川と河口の堆積システムを比較し、蛇行と網状の典型条件を区別する。

\paragraph{必要知識} 河川端成分、氾濫原とポイントバー、洪水時の級化、デルタ前進、相対海水準と供給、河口沈水。

\paragraph{考え方・正答を選ぶ理由（査読待ち初稿）} 五つは河川の端成分と河口堆積システムの対応を問う。蛇行の典型を、急勾配・大きな流量変動・粗粒子の大きな負荷で説明する肢を除く。この組合せは網状の典型に対応する。ただし実際の流路型には河岸の強度や供給も関わる。デルタは前進する堆積体で、海退と相対海水準低下は同義ではない。逆級化も洪水立上がりで形成し得るが、非一意の指標である。原典の典型モデルと、全称命題としての限界を分けて解く。本問の該当番号は3である。

\Needspace{4\baselineskip}\paragraph{選択肢の検討（査読待ち初稿）}\mbox{}\par

\noindent\textbf{1：適切（条件を含めて検討）} 河川を蛇行・網状の二端成分で考える。端成分モデルとして妥当である。中間形態や両特徴を持つ河川を除外する絶対二分ではない。\par

\noindent\textbf{2：適切（条件を含めて検討）} 氾濫原の砂岩に見られる逆級化を河川の特徴とする。洪水の増水過程では砂が上方粗粒化する場合がある。ただし逆級化は河川に固有ではなく、全氾濫原砂がそうなるという意味ではない。\par

\noindent\textbf{3：誤り・最も不適切} 蛇行河川を流量変動・勾配・粗粒負荷の大きい河川とする。これは網状河川の典型条件との混同である。蛇行は流路の曲率と側方移動、ポイントバーと泥質氾濫原で捉える。\par

\noindent\textbf{4：適切（条件を含めて検討）} 海湖側へ突出するデルタを海退期のシステムとする。デルタの前進が海退を示す場合の説明である。海退を海水準低下と同義にせず、デルタの存在だけで全時期を決めない。\par

\noindent\textbf{5：適切（条件を含めて検討）} 海進で河口が沈水する内湾をエスチュアリーとする。沈水谷を河川と潮汐が埋める河口内湾というモデルに対応する。海進する全河口の唯一の形ではない。\par

\paragraph{再利用する考え方} 観察できる構造と成因・環境の解釈を区別する。代表モデルに適用条件を添え、単一の形や測定値から唯一の環境を断定しない。

\paragraph{共通理論} \hyperref[theory:sediment:river]{第\ref*{theory:sediment:river}節}, \hyperref[theory:sediment:grading]{第\ref*{theory:sediment:grading}節}, \hyperref[theory:sediment:delta]{第\ref*{theory:sediment:delta}節}

\par 参考：\cite{sed_mit_environment,sed_walther_oregon}。

\paragraph{関連問題} \hyperref[q:h29-reg-applied-25]{H29-REG-APPLIED-25}

\par 難易度 3（校準初稿）、優先度未評価。

\section{堆積相の組合せ・累重とワルターの法則}

\subsection{令和6年度 問31}
\label{q:r06-reg-applied-31}

\noindent\textbf{原典：} sources/R06/17\_applied\_science.pdf、PDF物理10頁。

\textbf{公式正答：} 4（single）。

\paragraph{問題の要旨} 堆積相の記載、組合せ、累重の名前と、漸移的な上下関係の法則を区別する。

\Needspace{4\baselineskip}\paragraph{確認した条件・記述（独自要約）}\mbox{}\par

\noindent\textbf{判断方向：} 最も不適切・誤りを選ぶ\par

\noindent\textbf{原典照合：} 15原典画像で18問の問題幹、全5肢、否定、専門語、化学式、候補順を確認。\par

\noindent\textbf{適用条件：} 流れと構造の典型対応は一意成因を保証しない。用語、時間間隙、保存、試料処理の条件を個別に示す。\par

\paragraph{問われる核心} 堆積相の記載、組合せ、累重の名前と、漸移的な上下関係の法則を区別する。

\paragraph{必要知識} 観察属性と堆積相、association/succession、シーケンス境界、Waltherの連続・漸移条件。

\paragraph{考え方・正答を選ぶ理由（査読待ち初稿）} 堆積相を観察上の特徴で分け、組合せをアソシエーション、累重の順序をサクセションと呼ぶ。4は組合せをシーケンスとする用語の混同である。シーケンスは不整合又は対応する整合面で境界を定める地層群であり、単なる相の集合より条件が多い。ワルターの法則は漸移的・整合的な累重を扱うため、原典の条件を消して任意の上下接触へ使わない。

\Needspace{4\baselineskip}\paragraph{選択肢の検討（査読待ち初稿）}\mbox{}\par

\noindent\textbf{1：適切（条件を含めて検討）} 物理・生物・化学の特徴から堆積環境を復元する。観察記載から成因・環境を推定する解析の説明である。単一指標からの唯一解とはしない。\par

\noindent\textbf{2：適切（条件を含めて検討）} 側方隣接する環境の堆積物だけが漸移的に重なる。ワルターの法則の骨子である。原典の漸移的という条件を保持し、不整合を挟む全上下接触へ広げない。\par

\noindent\textbf{3：適切（条件を含めて検討）} 岩相・粒径・構造・化石で堆積相を区別する。相を区別する記載属性の組合せに対応する。環境の解釈と区別する。\par

\noindent\textbf{4：誤り・最も不適切} 堆積相の組合せを堆積シーケンスと呼ぶ。組合せの名称は堆積相アソシエーションである。シーケンスは境界で区分する成因的な地層群であり同じではない。\par

\noindent\textbf{5：適切（条件を含めて検討）} 累重・配列を堆積相サクセションと呼ぶ。相の上下の累重・配列というサクセションの説明に対応する。\par

\paragraph{再利用する考え方} 観察できる構造と成因・環境の解釈を区別する。代表モデルに適用条件を添え、単一の形や測定値から唯一の環境を断定しない。

\paragraph{共通理論} \hyperref[theory:sediment:facies]{第\ref*{theory:sediment:facies}節}, \hyperref[theory:sediment:walther]{第\ref*{theory:sediment:walther}節}

\par 参考：\cite{sed_mit_environment,sed_walther_oregon}。

\paragraph{関連問題} \hyperref[q:h23-reg-applied-25]{H23-REG-APPLIED-25}

\par 難易度 3（校準初稿）、優先度未評価。

\section{堆積環境指標と測定法の限界}

\subsection{平成23年度 問35}
\label{q:h23-reg-applied-35}

\noindent\textbf{原典：} sources/H23/17\_applied\_science.pdf、PDF物理11頁。

\textbf{公式正答：} 3（single）。

\par\textbf{検討未確定：} 陰イオン量の典型比較とECの古環境対応は濃度・試料処理・後の履歴に依存する。SEMの原典用語に対応する手法全体の追試は未確認。一次資料で確認した範囲と限界を明記して部分校正に残す。

\paragraph{問題の要旨} 堆積環境を反映する地球化学的指標と、元素分析・粒子表面観察の測定対象を分ける。

\Needspace{4\baselineskip}\paragraph{確認した条件・記述（独自要約）}\mbox{}\par

\noindent\textbf{判断方向：} 最も不適切・誤りを選ぶ\par

\noindent\textbf{原典照合：} 15原典画像で18問の問題幹、全5肢、否定、専門語、化学式、候補順を確認。\par

\noindent\textbf{適用条件：} 流れと構造の典型対応は一意成因を保証しない。用語、時間間隙、保存、試料処理の条件を個別に示す。\par

\paragraph{問われる核心} 堆積環境を反映する地球化学的指標と、元素分析・粒子表面観察の測定対象を分ける。

\paragraph{必要知識} XRFの特性X線と元素量、SEMの表面観察、硫化物と酸化還元、ECの試料処理依存。

\paragraph{考え方・正答を選ぶ理由（査読待ち初稿）} 蛍光X線分析は試料の原子を励起し、その特性X線を使う元素分析である。内部の葉理や粒子配列を観察する方法という3が誤りである。黄鉄鉱 \(\mathrm{FeS_2}\) は環境に応じて生成・酸化し、粘土の懸濁液のECもその反応に関わる。GSJの2011年資料は試料処理によるECの変化と環境推定の限界を示す。SEMの表面観察では風による運搬を推定する証拠を得られるが、後の化学的改変も検討する。

\Needspace{4\baselineskip}\paragraph{選択肢の検討（査読待ち初稿）}\mbox{}\par

\noindent\textbf{1：適切（条件を含めて検討）} FeS2含有量は堆積環境に依存する。硫化物形成は硫黄・鉄の供給、還元状態などに依存する。酸化・溶脱後の値を堆積時と無条件に同一視しない。\par

\noindent\textbf{2：適切（条件を含めて検討）} 自然水の陰イオン量は淡水より海水が多い。通常の海水は塩化物等の溶存イオン濃度が淡水より大きい。塩湖等の例外を含む全ての水への全称ではない。\par

\noindent\textbf{3：誤り・最も不適切} XRFで内部の堆積構造を観察する。蛍光X線のエネルギーと強度から元素を分析するので、内部堆積構造の画像を直接得る方法ではない。\par

\noindent\textbf{4：適切（条件を含めて検討）} 石英表面のSEMは風成作用の推定に有効。粒子表面の微細摩耗・衝突形態を環境推定に使える。ただし水成・風成・溶食や再運搬の履歴を単独で一意決定しない。\par

\noindent\textbf{5：適切（条件を含めて検討）} 粘土混濁水の電気伝導度変化は古環境に対応する。懸濁液のECを古環境指標とする研究があるが、乾燥と酸化、混合比や温度、後の溶脱にも依存し、単独判定の限界がある。\par

\paragraph{再利用する考え方} 観察できる構造と成因・環境の解釈を区別する。代表モデルに適用条件を添え、単一の形や測定値から唯一の環境を断定しない。

\paragraph{共通理論} \hyperref[theory:sediment:proxy]{第\ref*{theory:sediment:proxy}節}

\par 参考：\cite{sed_xrf_jaea,sed_gsj_ec2011,sed_sem_shimane}。

\par 難易度 3（校準初稿）、優先度未評価。

\section{未固結堆積物の二次変形}

\subsection{令和元年度 問28}
\label{q:r01-reg-applied-28}

\noindent\textbf{原典：} sources/R01/17\_applied\_science.pdf、PDF物理10頁。

\textbf{公式正答：} 2（single）。

\paragraph{問題の要旨} 堆積中にできる構造と、その後の未固結段階に変形してできる構造を区別する。

\Needspace{4\baselineskip}\paragraph{確認した条件・記述（独自要約）}\mbox{}\par

\noindent\textbf{判断方向：} 堆積後の未固結段階の二次構造を選ぶ\par

\noindent\textbf{原典照合：} 15原典画像で18問の問題幹、全5肢、否定、専門語、化学式、候補順を確認。\par

\noindent\textbf{適用条件：} 流れと構造の典型対応は一意成因を保証しない。用語、時間間隙、保存、試料処理の条件を個別に示す。\par

\paragraph{問われる核心} 堆積中にできる構造と、その後の未固結段階に変形してできる構造を区別する。

\paragraph{必要知識} 初生/二次の分類軸、未固結変形、荷重差・間隙水、葉理と後の変形の区別。

\paragraph{考え方・正答を選ぶ理由（査読待ち初稿）} 分類の軸は形の大きさではなく形成時期と過程である。流痕・リップル・葉理・級化は流れと沈積による初生構造に対応する。火炎構造は未固結の泥が上の砂へ突き出す荷重・流動化に関わるため二次構造であり、本問では2である。後から曲げられた葉理を見た場合は、葉理の形成と変形を別の段階として読む。火炎があるだけで地震と断定しない。

\Needspace{4\baselineskip}\paragraph{選択肢の検討（査読待ち初稿）}\mbox{}\par

\noindent\textbf{1：非該当・不適切} 流痕。流れが堆積面を侵食する初生構造である。\par

\noindent\textbf{2：該当・適切} 火炎構造。砂泥の荷重差や流動化による、堆積後の未固結変形の二次構造である。\par

\noindent\textbf{3：非該当・不適切} リップル。堆積しながら流れが床形を作る初生構造である。\par

\noindent\textbf{4：非該当・不適切} 葉理。堆積粒子の層状配列を記録する初生構造である。後から変形した葉理という別概念と区別する。\par

\noindent\textbf{5：非該当・不適切} 級化層理。堆積中の粒径の時間変化を記録する初生構造である。\par

\paragraph{再利用する考え方} 観察できる構造と成因・環境の解釈を区別する。代表モデルに適用条件を添え、単一の形や測定値から唯一の環境を断定しない。

\paragraph{共通理論} \hyperref[theory:sediment:soft]{第\ref*{theory:sediment:soft}節}, \hyperref[theory:sediment:cross]{第\ref*{theory:sediment:cross}節}, \hyperref[theory:sediment:grading]{第\ref*{theory:sediment:grading}節}

\par 参考：\cite{sed_soft_deformation,sed_houston_structures}。

\paragraph{関連問題} \hyperref[q:r02-reg-applied-29]{R02-REG-APPLIED-29}, \hyperref[q:r06-reg-applied-30]{R06-REG-APPLIED-30}

\par 難易度 3（校準初稿）、優先度未評価。

\subsection{令和6年度 問30}
\label{q:r06-reg-applied-30}

\noindent\textbf{原典：} sources/R06/17\_applied\_science.pdf、PDF物理10頁。

\textbf{公式正答：} 3（single）。

\paragraph{問題の要旨} 堆積中にできる構造と、その後の未固結段階に変形してできる構造を区別する。

\Needspace{4\baselineskip}\paragraph{確認した条件・記述（独自要約）}\mbox{}\par

\noindent\textbf{判断方向：} 堆積後の未固結段階の二次構造を選ぶ\par

\noindent\textbf{原典照合：} 15原典画像で18問の問題幹、全5肢、否定、専門語、化学式、候補順を確認。\par

\noindent\textbf{適用条件：} 流れと構造の典型対応は一意成因を保証しない。用語、時間間隙、保存、試料処理の条件を個別に示す。\par

\paragraph{問われる核心} 堆積中にできる構造と、その後の未固結段階に変形してできる構造を区別する。

\paragraph{必要知識} 初生/二次の分類軸、未固結変形、荷重差・間隙水、葉理と後の変形の区別。

\paragraph{考え方・正答を選ぶ理由（査読待ち初稿）} 分類の軸は形の大きさではなく形成時期と過程である。流痕・リップル・葉理・級化は流れと沈積による初生構造に対応する。火炎構造は未固結の泥が上の砂へ突き出す荷重・流動化に関わるため二次構造であり、本問では3である。後から曲げられた葉理を見た場合は、葉理の形成と変形を別の段階として読む。火炎があるだけで地震と断定しない。

\Needspace{4\baselineskip}\paragraph{選択肢の検討（査読待ち初稿）}\mbox{}\par

\noindent\textbf{1：非該当・不適切} 葉理。堆積粒子の層状配列を記録する初生構造である。後から変形した葉理という別概念と区別する。\par

\noindent\textbf{2：非該当・不適切} 級化層理。堆積中の粒径の時間変化を記録する初生構造である。\par

\noindent\textbf{3：該当・適切} 火炎構造。砂泥の荷重差や流動化による、堆積後の未固結変形の二次構造である。\par

\noindent\textbf{4：非該当・不適切} リップル。堆積しながら流れが床形を作る初生構造である。\par

\noindent\textbf{5：非該当・不適切} 流痕。流れが堆積面を侵食する初生構造である。\par

\paragraph{再利用する考え方} 観察できる構造と成因・環境の解釈を区別する。代表モデルに適用条件を添え、単一の形や測定値から唯一の環境を断定しない。

\paragraph{共通理論} \hyperref[theory:sediment:soft]{第\ref*{theory:sediment:soft}節}, \hyperref[theory:sediment:cross]{第\ref*{theory:sediment:cross}節}, \hyperref[theory:sediment:grading]{第\ref*{theory:sediment:grading}節}

\par 参考：\cite{sed_soft_deformation,sed_houston_structures}。

\paragraph{関連問題} \hyperref[q:r01-reg-applied-28]{R01-REG-APPLIED-28}

\par 難易度 3（校準初稿）、優先度未評価。

\section{津波堆積物の形成・認定・保存}

\subsection{令和3年度 問30}
\label{q:r03-reg-applied-30}

\noindent\textbf{原典：} sources/R03/17\_applied\_science.pdf、PDF物理11頁。

\textbf{公式正答：} 2（single）。

\par\textbf{検討未確定：} 公式/独立2を保持するが、単純/複雑の頻度は環境と尺度に依存し、USGSには構造が稀な例も多い。②の代表頻度を厳密には未確定。④の火山起源津波堆積物の個別報告も未確認。①③⑤の一般化の限界を含めAstra重点査読へ残す。

\paragraph{問題の要旨} 津波堆積物を、供給源、反復する流れ、非地震起源と形成後の保存条件から検討する。

\Needspace{4\baselineskip}\paragraph{確認した条件・記述（独自要約）}\mbox{}\par

\noindent\textbf{判断方向：} 最も不適切・誤りを選ぶ\par

\noindent\textbf{原典照合：} 15原典画像で18問の問題幹、全5肢、否定、専門語、化学式、候補順を確認。\par

\noindent\textbf{適用条件：} 流れと構造の典型対応は一意成因を保証しない。用語、時間間隙、保存、試料処理の条件を個別に示す。\par

\paragraph{問われる核心} 津波堆積物を、供給源、反復する流れ、非地震起源と形成後の保存条件から検討する。

\paragraph{必要知識} 反復波と遡上/戻り流れ、供給源、微化石、背景堆積と保存、非地震起源、事例頻度の母集団。

\paragraph{考え方・正答を選ぶ理由（査読待ち初稿）} 初回は波の反復・戻り流れによる多様性から2を最も不適切と判定し、公式2に一致した。藤野の報告は塊状から複数級化ユニットまでの多様性を示す。一方、USGSの近年の事例集では堆積構造が稀な例も多い。したがって「津波なら必ず複雑だから2が誤り」とは書けず、対象環境・記録の尺度による頻度の違いを疑義に残す。生物遺骸、分布、構造、年代などの複数証拠で認定し、形成と保存と観察可能性を分ける。

\Needspace{4\baselineskip}\paragraph{選択肢の検討（査読待ち初稿）}\mbox{}\par

\noindent\textbf{1：適切（条件を含めて検討）} 生物遺骸から侵食した海底の深度を検討する。生物の生息深度は供給源を制約する証拠となるが、混合と再堆積の影響を考慮する。\par

\noindent\textbf{2：誤り・最も不適切} 堆積構造は比較的単純な場合が多い。多様な層理・正逆級化・複数波のユニットがあるため一律な単純化が問題となる。ただしUSGS近年資料では構造が稀な例も多く、単純/複雑の一般頻度は資料依存である。\par

\noindent\textbf{3：適切（条件を含めて検討）} 礫床河川・扇状地・岩礁海岸は調査に不向き。通常の粗粒供給や侵食でイベント層の識別・保存が難しい。岩礁に津波記録が絶対にないという意味ではない。\par

\noindent\textbf{4：適切（条件を含めて検討）} 火山・隕石・海底地すべり起源の津波にも報告がある。隕石衝突とストレッガーの報告は確認した。火山起源堆積物の個別報告は今回未確認であり断定的な全確認扱いにしない。\par

\noindent\textbf{5：適切（条件を含めて検討）} 内湾は安定した環境で保存されやすい。静かな内湾・湿地の背景堆積と後の埋積は保存に有利な場合がある。侵食や擾乱を受けない保証ではない。\par

\paragraph{再利用する考え方} 観察できる構造と成因・環境の解釈を区別する。代表モデルに適用条件を添え、単一の形や測定値から唯一の環境を断定しない。

\paragraph{共通理論} \hyperref[theory:sediment:tsunami]{第\ref*{theory:sediment:tsunami}節}

\par 参考：\cite{sed_tsunami_fujino,sed_tsunami_usgs,sed_tsunami_2002,sed_tsunami_gsj_survey}。

\par 難易度 4（校準初稿）、優先度未評価。

\section{形成時の上下と堆積構造}

\subsection{令和2年度 問29}
\label{q:r02-reg-applied-29}

\noindent\textbf{原典：} sources/R02/17\_applied\_science.pdf、PDF物理10頁。

\textbf{公式正答：} 1（single）。

\paragraph{問題の要旨} 転倒した地層の形成時の上下を決める構造と、節理の伝播を示す構造を区別する。

\Needspace{4\baselineskip}\paragraph{確認した条件・記述（独自要約）}\mbox{}\par

\noindent\textbf{判断方向：} 最も不適切・誤りを選ぶ\par

\noindent\textbf{原典照合：} 15原典画像で18問の問題幹、全5肢、否定、専門語、化学式、候補順を確認。\par

\noindent\textbf{適用条件：} 流れと構造の典型対応は一意成因を保証しない。用語、時間間隙、保存、試料処理の条件を個別に示す。\par

\paragraph{問われる核心} 転倒した地層の形成時の上下を決める構造と、節理の伝播を示す構造を区別する。

\paragraph{必要知識} 節理と堆積構造、脱水皿の幾何、荷重方向、侵食型と底面型、生痕の開口。

\paragraph{考え方・正答を選ぶ理由（査読待ち初稿）} 羽毛状構造は固結岩の割れ目が広がった過程を記録する。堆積時の上を示す幾何ではないため1が最も不適切である。皿状・荷重・流痕は形成時の重力や層の界面との対応から上下を考え、生痕は開口や埋積を含めて読む。現在の鉛直と堆積時の鉛直は転倒で違うため、一つの曖昧な形だけでなく複数構造を合わせる。

\Needspace{4\baselineskip}\paragraph{選択肢の検討（査読待ち初稿）}\mbox{}\par

\noindent\textbf{1：誤り・最も不適切} 羽毛状構造。節理の破断面の起点から広がる伝播模様である。堆積時の上下面を直接示さない。\par

\noindent\textbf{2：適切（条件を含めて検討）} 皿状構造。一般に上に凹の脱水構造の形と周囲の層を使い上下を検討できる。\par

\noindent\textbf{3：適切（条件を含めて検討）} 荷重痕。上の砂が下の泥へ突出するため、荷重の向きを使える。\par

\noindent\textbf{4：適切（条件を含めて検討）} 流痕。侵食した凹面と上位層底面の凸型の充填を区別すると上下を検討できる。\par

\noindent\textbf{5：適切（条件を含めて検討）} 生痕化石。巣穴の開口・分岐や地表からの貫入などが明瞭なら上下を使える。全ての名称だけで上下が分かるとは限らない。\par

\paragraph{再利用する考え方} 観察できる構造と成因・環境の解釈を区別する。代表モデルに適用条件を添え、単一の形や測定値から唯一の環境を断定しない。

\paragraph{共通理論} \hyperref[theory:sediment:wayup]{第\ref*{theory:sediment:wayup}節}, \hyperref[theory:sediment:soft]{第\ref*{theory:sediment:soft}節}

\par 参考：\cite{sed_plumose_michigan,sed_mit_environment,sed_houston_structures}。

\paragraph{関連問題} \hyperref[q:r01-reg-applied-28]{R01-REG-APPLIED-28}

\par 難易度 3（校準初稿）、優先度未評価。

